\section{Future Research Boundary}

The negative results motivate questions about task-specific symmetry,
representation, and residual structure, but the completed evidence does not
answer them. A future geometric-QML direction was discussed as a hypothesis,
not as a result. Here a symmetry would be a transformation, such as rotating
all relevant vectors together, that preserves the physical relationship between
inputs and labels. An equivariant model is designed so that its output changes
consistently under such a transformation. The first unanswered question is
whether any nontrivial symmetry actually survives fixed Earth--Moon--Sun
ephemerides, thrust constraints, and reference-frame choices. Equivariance
alone would not imply a QML advantage over classical equivariant models.

A future protocol would require a mathematical contract, a symmetry audit,
baseline reproduction, a representation- and parameter-matched classical
comparison, ablations of symmetry and initialization, resource accounting,
and explicit rejection criteria. The D054-A/P018 binding freeze records how a
future audit could specify its representation, dynamical Lie algebra (DLA)
basis, allowed circuit generators, development-only case groups, numerical
tolerance, output paths, and figures. A DLA describes the transformations that
can be generated by the allowed circuit operations. Closure decision D055-C
declined to execute the audit. Accordingly, this manuscript contains no
DOP853 trajectory replay (a high-order numerical integration check), no
singular-value-decomposition (SVD) rank or nullspace result, no new symmetry
dataset, no variational circuit training, and no geometric-QML performance
claim.

This boundary is not a weakness in reporting. It prevents a promising idea
from being laundered into evidence through narrative. The appropriate next
step, if future authors choose to reopen research, is a separately approved
protocol that can reject the hypothesis as readily as it can support it.
